\section{Research design: a Design Science framing}
\label{sec:design}

\subsection{Problem class and relevance}

The problem class is organisational rather than computational. Organisations are required, with
increasing regularity, to demonstrate that a system behaves as a specified rule intends --- for audit,
for regulatory compliance, for contractual acceptance, or simply for operational consistency. BDD is
one of the instruments adopted to that end. But an organisation that adopts BDD acquires a
\emph{representation} of its rules, not a guarantee about them, and the representation can retain its
appearance while losing its function.

The relevance of the problem class is a matter of record in the decay literature: disabled tests,
rotten green tests and inadequate assertions are normal, not pathological
\cite{Kim2021Disabled,Delplanque2019RottenGreen,zhi2013inadequateassertions}. The specificity of the
BDD case is the appearance/function asymmetry described in Section~\ref{sec:background}: a
\texttt{.feature} file reads like a specification to a business stakeholder even when its step
definitions stopped checking anything. A manager reviewing the artefact sees compliance; the
instrument we propose tests whether compliance is real.

\subsection{The artefact}

Following the design-science convention of describing an artefact's rationale and evaluation rather
than only its implementation, we state the artefact as a protocol with five steps:

\begin{enumerate}
\item \textbf{Delimit the domain layer} using a three-tier rule (path convention, framework
      exclusion, decision-bearing fallback), so that mutation is applied where business decisions
      live rather than where data is moved (Section~\ref{sec:method}, protocol step 1).
\item \textbf{Measure oracle strength twice}, in BDD-only and full-suite modes, with coverage
      conditioning mandatory and with the differential attributed to the specification layer
      explicitly (protocol step 2).
\item \textbf{Compute seven drift indicators} with fixed polarity and operational definitions, two of
      them primary and five exploratory (protocol step 3).
\item \textbf{Test for accumulation over history} using a deterministic sample of commits and a
      robust trend statistic, plus a two-snapshot change estimate on a \emph{common} mutant set
      (protocol step 4).
\item \textbf{Interpret under rules that forbid the common misreadings} --- most importantly, that
      ``not covered'' is not ``drifted'', that a snapshot is not a trend, and that no association
      below the detectable-effect floor should be presented as evidence of absence (protocol step 5).
\end{enumerate}

The full protocol is a deliverable of this work (\texttt{docs/PROTOCOL.md}) and is written for
practitioners, not for reviewers. Three design decisions in it deserve explicit defence because they
each foreclose a convenient shortcut:

\begin{description}
\item[Coverage conditioning is mandatory, not advisory.] Given that the raw coverage--mutation
      correlation largely vanishes once analysis is restricted to covered code
      \cite{Jain2023MindTheGap}, a study that reports only an overall mutation score is measuring
      coverage under another name. We therefore require both values side by side in every table.
\item[Differential attribution is required.] A suite can be strong for reasons unrelated to BDD:
      a mature project with excellent unit tests will have a high overall mutation score whether or
      not its Gherkin layer checks anything. Reporting only the total would confound BDD with general
      engineering quality. The differential $\dk$ on a common mutant set is the artefact's answer to
      this, and the paper states plainly that no causal claim about BDD is available from an
      observational design.
\item[Stratification by tool is required.] Because mutation tools are not interchangeable
      \cite{Kintis2018HowEffectiveTools}, pooling scores across ecosystems would manufacture a
      common scale that does not exist. The artefact therefore forbids a single pooled correlation as
      the headline statistic.
\end{description}

\subsection{Evaluation plan}

The artefact is evaluated on two dimensions, and we are explicit about which succeeded and which did
not:

\begin{itemize}
\item \textbf{Constructed and tested (feasibility).} Can the measurement be performed at all with
      available, pinned tooling, in more than one ecosystem, with the per-mutant coverage and
      per-scenario attribution the artefact requires? This is answered by construction: we built and
      ran the harnesses, and we report the observed mutation counts and scores
      (Section~\ref{sec:results}).
\item \textbf{Applied to a population (attrition and observed values).} What happens when the
      instrument is pointed at a defined frame of public repositories under pre-registered
      eligibility rules? This is answered by observation, and the answer includes how many candidate
      repositories could not be measured and why. Section~\ref{sec:results} reports the funnel; the
      interpretation of an instrument whose population is thin is taken up in
      Section~\ref{sec:discussion}.
\end{itemize}

\subsection{Relation between questions and observations}

Table~\ref{tab:rq-map} makes explicit which observation answers which question, and what the answer
rests on. This mapping was fixed before measurement; making it explicit is what prevents a paper from
retrospectively promoting an exploratory finding to a confirmatory one.

\begin{table}[htbp]
\centering
\caption{Research questions, the observations intended to answer them, and the evidentiary basis of
each. The third column names what could make each answer fail, which is where the threats in
Section~\ref{sec:threats} attach.}
\label{tab:rq-map}
\begin{tblr}{width=\textwidth,colspec={X[1.0]X[1.6]X[1.6]},rowsep=2pt}
\toprule
\textbf{RQ} & \textbf{Observation} & \textbf{What the answer depends on} \\
\midrule
RQ1 & Two-rater agreement on the judgement-coded indicators; seeded-drift detection rate;
feasibility of the mutation harness in each ecosystem & Whether independent human raters are
available; whether injected drift is detectable in the direction the indicator claims \\
RQ2 & $\mscov$ over domain mutants in BDD-only mode; $\dk$ against the full suite on a common
mutant set & Whether a green baseline and an identifiable domain layer exist in the repository; tool
comparability within a stratum \\
RQ3 & Association between pre-specified indicators and $\mscov$ within strata, with bootstrap CIs;
Theil--Sen slopes across sampled commits & Statistical power at the achieved $n$; whether the history
contains enough relevant commits to support a slope \\
\bottomrule
\end{tblr}
\end{table}

\subsection{Ethics and governance}

All analysis is confined to public repositories under permissive licences, with the licence recorded
per repository and repositories without a detectable permissive licence excluded before any
modification. Mutation analysis necessarily modifies code; it is performed only in isolated working
copies, never against a repository in place, and seeded-drift experiments are confined to separate
copies so that no measured artefact is altered. No human participants are involved, so no ethics
committee approval applies; the research uses only publicly available artefacts.
